\clearpage

\section{Discussion}

Uncertain early exponential growth rate calculations feed into estimates of $R_0$ --- they are a critical ingredient for most $R_0$ estimation formulae. 

The superspreading mechanism introduced here is novel, and implies that some amount of superspreading cannot be traced {\it a priori} to the individual: there is no such thing as identifying superspreaders, since the superspreader in this case may not have any more inherent infectiousness or a greater number of contacts; instead, they were just infected at the wrong time. 

Accounting for the individual infectiousness profile is critical for anticipating the imapct of outbreak interventions. Simply calculating reductions in infectiousness is not enough: we get differences in overdispersion and the generation interval distribution that are sensitive to the punctuatedness of $a_i(\tau)$ and can lead to profoundly different epidemic dynamcis than a naive approach of simulating an equivalent epidemic without changing overdispersion or the generation interval distribution. 

Estiamting the shape of the individual infectiousness profile can be challenging given the need for precise estimates of time of infection and a good balance of depth and breadth of sampling. Fortunately, identifiability of the shape of $a_i(\tau)$ is better in situations where it matters most, \ie, when $R_0$ is high. 

{\bf Limitations.}

{\bf Concluding thoughts.}